\section{Confusion matrices and a repeat-run check}
\label{app:confusion}

\subsection{Full confusion-matrix set}
\label{app:confusion:matrices}

Table~\ref{tab:confmat} gives the complete confusion-matrix set underlying the per-class analysis of \ref{sec:results:perclass}: 5 arms $\times$ 3 conditions (clean, $0$\,dB, $-6$\,dB), each summed over 5 seeds (support OR$=17{,}520$, IR$=8{,}720$ per cell). These counts were dumped per-sample and re-verified against the reported recall/precision/F1 metrics cell-by-cell before use.

\onecolumn\begin{longtable}{@{}llrrrr@{}}
\caption{Full confusion-matrix set (summed over 5 seeds; support OR=17{,}520, IR=8{,}720 per cell). Columns give true$\to$predicted counts.}\label{tab:confmat}\\
\toprule
Arm & Condition & OR$\to$OR & OR$\to$IR & IR$\to$OR & IR$\to$IR \\
\midrule
\endfirsthead
\toprule
Arm & Condition & OR$\to$OR & OR$\to$IR & IR$\to$OR & IR$\to$IR \\
\midrule
\endhead
\bottomrule
\endfoot
BM3 & clean & 17{,}500 & 20 & 3{,}948 & 4{,}772 \\
 & 0\,dB & 17{,}395 & 125 & 1{,}364 & 7{,}356 \\
 & $-6$\,dB & 16{,}613 & 907 & 1{,}664 & 7{,}056 \\
BM3-frozen & clean & 17{,}303 & 217 & 3{,}917 & 4{,}803 \\
 & 0\,dB & 17{,}089 & 431 & 2{,}060 & 6{,}660 \\
 & $-6$\,dB & 15{,}634 & 1{,}886 & 2{,}543 & 6{,}177 \\
S4D+gate & clean & 17{,}327 & 193 & 4{,}969 & 3{,}751 \\
 & 0\,dB & 16{,}303 & 1{,}217 & 3{,}322 & 5{,}398 \\
 & $-6$\,dB & 13{,}267 & 4{,}253 & 3{,}244 & 5{,}476 \\
frozen$-$gate & clean & 16{,}935 & 585 & 5{,}139 & 3{,}581 \\
 & 0\,dB & 16{,}092 & 1{,}428 & 5{,}406 & 3{,}314 \\
 & $-6$\,dB & 13{,}200 & 4{,}320 & 5{,}469 & 3{,}251 \\
S4D & clean & 16{,}828 & 692 & 3{,}354 & 5{,}366 \\
 & 0\,dB & 14{,}945 & 2{,}575 & 2{,}897 & 5{,}823 \\
 & $-6$\,dB & 12{,}477 & 5{,}043 & 3{,}810 & 4{,}910 \\
\end{longtable}
\twocolumn

\subsection{Repeat-run check}
\label{app:confusion:floor}

Because two independent process runs of an identical configuration are not bitwise reproducible on this hardware (non-deterministic GPU kernel reductions in the selective-scan and normalisation ops), we measured a process-level noise floor directly: the same arm, condition, and seed (\texttt{bm3\_frozen}, clean, seed 0, 2 epochs) was run twice in separate processes with identical evaluation fingerprints and identical parameter counts confirmed. The two runs' best macro-F1 differed by
\[
|0.69025 - 0.69712| = 0.00687 \approx 0.687\pp,
\]
which we report as an observed repeat difference of $\sim\!0.7\pp$ (cited in Section~\ref{sec:protocol:prereg}). It is a single repeat of one arm, one condition and a 2-epoch run, and therefore neither an estimate nor a bound of process variability for full-length runs. It is an order of magnitude below the seed-level variability observed within the deep-noise band across the curve tables underlying Figures~\ref{fig:snr_curves} and~\ref{fig:graft} (standard deviations typically $2$--$9\pp$, e.g.\ BM3-frozen $5.7$--$9.4\pp$ and S4D $1.7$--$4.5\pp$ at $\{0,-2,-6\}$\,dB). Arm-level effects in Section~\ref{sec:results} are judged against seed-level variability, not against this single repeat.
